% 英文摘要

Marine propellers are subjected to non-uniform wake, blade-frequency excitation, and fluid–structure interaction during operation, which may induce complex vibration responses. Their modal characteristics are directly related to the structural safety, vibration and noise performance, and service reliability of propulsion systems. Although conventional finite element methods can calculate propeller modal frequencies with relatively high accuracy, they require repeated geometric modeling, meshing, and finite element solution in multi-parameter design and large-scale scheme evaluation. As a result, their computational cost is high and their efficiency is insufficient for rapid prediction. To address this problem, this study proposes a rapid prediction method for the wet modal frequencies of marine propellers by integrating high-fidelity finite element simulation, modal test validation, and machine-learning surrogate modeling.

With respect to high-fidelity finite element simulation, a computational procedure for dry and wet modal analysis of propellers is established based on modal analysis theory and finite element theory. Key aspects such as element type selection, mesh size determination, material parameter setting and calibration, and coupled structural–fluid-domain modeling for wet modal analysis are specified. Taking the DTMB4382 propeller as an example, the first 20 dry and wet modal frequencies and the corresponding added-water frequency correction coefficients are obtained, providing a basis for subsequent rapid prediction of wet modal frequencies. In terms of database construction, parametric modeling and automated simulation techniques are combined. Three-dimensional propeller models are generated using Python and CadQuery, automatic meshing is performed through HyperMesh batch scripts, and modal calculations and frequency extraction are implemented using Abaqus Python scripts. A dry modal frequency database is thereby constructed to support machine-learning model training. In terms of modal testing, a dry/wet modal test procedure based on the roving hammer method, acceleration-response frequency response function measurement, and modal parameter identification is developed. Modal tests in air and water are conducted on a seven-bladed propeller and a ten-bladed propeller. The comparison between simulation and test results for the seven-bladed propeller verifies the accuracy and engineering applicability of the proposed simulation procedure.

For machine-learning prediction, dry modal frequency prediction models are developed based on the constructed dry modal frequency database. Pearson and Spearman correlation coefficients are first used to analyze the relationships between the original geometric parameters and the first 20 dry modal frequencies. A total of 24 key geometric parameters, including propeller diameter, blade number, and the chord length and thickness at 11 radial sections, are selected as model inputs, thus reducing the dimensionality of the input features. Four machine-learning models, namely Random Forest, Gradient Boosting Regression, Gaussian Process Regression, and Multi-Layer Perceptron, are then constructed to predict the first 20 dry modal frequencies of the propeller. Random search is employed for hyperparameter optimization. The prediction performance of the models is evaluated using mean absolute error (MAE), root mean square error (RMSE), coefficient of determination (R²), and mean absolute percentage error (MAPE). The results show that the Multi-Layer Perceptron model achieves the best overall performance, with an overall MAE of 5.133937, RMSE of 7.445019, R² of 0.996914, and MAPE of 1.889931\% on the test set. This indicates that the model can accurately learn the nonlinear mapping between key propeller geometric parameters and dry modal frequencies. On this basis, the predicted first 20 dry modal frequencies obtained from the Multi-Layer Perceptron model are combined with the added-water frequency correction coefficients to establish a rapid prediction process for wet modal frequencies using propeller geometric parameters as inputs.

The results demonstrate that the proposed method can significantly reduce the repeated finite element computational cost in propeller modal frequency evaluation while maintaining satisfactory prediction accuracy. It provides a feasible technical approach for modal frequency avoidance analysis, structural optimization design, and rapid wet modal frequency prediction of marine propellers.